\chapter{\texorpdfstring{\(SU(2)\) to \(SO(3)\)}{SU(2) to SO(3)}}
\label{ch:viii13-su2-so3}

The double cover \(SU(2)\to SO(3)\) is a central example where topology,
group structure, and geometry reinforce one another.  The group \(SU(2)\) is
a 3-sphere, while \(SO(3)\) is obtained by identifying antipodal points.  Unit
quaternions provide the most transparent formula for the covering map.

\section{The group \(SU(2)\)}

\begin{definition}[Special unitary group]
\label{def:viii13-su2}
\[
SU(2)=\{A\in M_2(\mathbb C):A^*A=I,\ \det A=1\}.
\]
\end{definition}

Every element has the form
\[
A=
\begin{pmatrix}
\alpha&-\overline\beta\\
\beta&\overline\alpha
\end{pmatrix},
\qquad
|\alpha|^2+|\beta|^2=1.
\]

\begin{proposition}[Three-sphere model]
\label{prop:viii13-s3}
As a topological space,
\[
SU(2)\cong S^3.
\]
\end{proposition}

Indeed \((\alpha,\beta)\in\mathbb C^2\cong\mathbb R^4\) lies on the unit
3-sphere.

\section{Unit quaternions}

Write
\[
\mathbb H=\{a+bi+cj+dk:a,b,c,d\in\mathbb R\}.
\]
The unit quaternions
\[
S^3=\{q\in\mathbb H:|q|=1\}
\]
form a group under quaternion multiplication and identify with \(SU(2)\).

The purely imaginary quaternions
\[
\operatorname{Im}\mathbb H=\{bi+cj+dk\}
\]
are naturally \(\mathbb R^3\).

\section{Quaternion conjugation}

For a unit quaternion \(q\), define
\[
R_q(v)=qvq^{-1},
\qquad v\in\operatorname{Im}\mathbb H.
\]

\begin{theorem}[Quaternion rotation]
\label{thm:viii13-quaternion-rotation}
\(R_q\) is an orientation-preserving orthogonal transformation of
\(\mathbb R^3\).  Hence
\[
R_q\in SO(3).
\]
\end{theorem}

Thus
\[
\rho:S^3\cong SU(2)\longrightarrow SO(3),
\qquad
q\longmapsto R_q
\]
is a group homomorphism.

\section{Kernel}

\begin{theorem}[Kernel of the rotation map]
\label{thm:viii13-kernel}
\[
\ker\rho=\{\pm1\}.
\]
\end{theorem}

\begin{proof}
Both \(1\) and \(-1\) commute with every imaginary quaternion and therefore
act trivially.  Conversely, a unit quaternion acting trivially by conjugation
commutes with all imaginary quaternions, hence lies in the real center of
\(\mathbb H\), so it is \(\pm1\).
\end{proof}

\section{Axis-angle formula}

Let \(u\in\operatorname{Im}\mathbb H\) be a unit vector.  Put
\[
q=\cos\frac{\theta}{2}+u\sin\frac{\theta}{2}.
\]

\begin{theorem}[Half-angle formula]
\label{thm:viii13-half-angle}
Conjugation by \(q\) rotates \(\mathbb R^3\) through angle \(\theta\) about the
axis \(u\).
\end{theorem}

The half-angle is the geometric source of the double cover.

\section{Surjectivity}

Every rotation in \(SO(3)\) has an axis and angle.  The half-angle formula
therefore gives:

\begin{theorem}[Surjectivity]
\label{thm:viii13-surjective}
The map
\[
\rho:SU(2)\to SO(3)
\]
is surjective.
\end{theorem}

Together with the kernel computation,
\[
SO(3)\cong SU(2)/\{\pm I\}.
\]

\section{Double covering}

\begin{theorem}[Double cover]
\label{thm:viii13-double-cover}
\[
\rho:SU(2)\to SO(3)
\]
is a two-sheeted covering map with deck group
\[
\{\pm I\}\cong\mathbb Z/2.
\]
\end{theorem}

The nontrivial deck transformation is multiplication by \(-I\).

\section{\(SO(3)\) as projective 3-space}

Since
\[
SU(2)\cong S^3
\]
and the kernel identifies \(q\) with \(-q\),
\[
SO(3)\cong S^3/\{\pm1\}\cong\mathbb RP^3.
\]

\begin{corollary}
\label{cor:viii13-rp3}
\[
SO(3)\cong\mathbb RP^3
\]
as topological manifolds.
\end{corollary}

\section{Fundamental group of \(SO(3)\)}

Because \(S^3\) is simply connected, the double cover is universal.

\begin{theorem}[Fundamental group]
\label{thm:viii13-pi1}
\[
\pi_1(SO(3))\cong\mathbb Z/2.
\]
\end{theorem}

\section{A \(2\pi\) rotation loop}

A path in \(SU(2)\) from \(1\) to \(-1\) projects to a loop in \(SO(3)\).
For example
\[
q(t)=\cos(\pi t)+u\sin(\pi t)
\]
projects to rotation through angle \(2\pi t\) about \(u\).

Its lift is not closed, so the projected loop represents the nontrivial
element of \(\pi_1(SO(3))\).

\section{A \(4\pi\) rotation}

Traversing twice gives a lift from \(1\) to \(1\).  Hence the \(4\pi\)-rotation
loop is null-homotopic in \(SO(3)\).

This is the topological content of the familiar distinction between \(2\pi\)
and \(4\pi\) rotation paths.

\section{Deck symmetry and lifting}

A loop in \(SO(3)\) lifts from \(I\in SU(2)\) either back to \(I\) or to
\(-I\).  These two possible endpoints detect the two elements of
\[
\pi_1(SO(3))\cong\mathbb Z/2.
\]

\section{Matrix picture}

The quaternion formula is equivalent to the adjoint action of \(SU(2)\) on its
three-dimensional Lie algebra.  After identifying the Lie algebra with
\(\mathbb R^3\), the adjoint representation is precisely
\[
\operatorname{Ad}:SU(2)\to SO(3).
\]

\section{Central quotient principle}

The covering map is a group quotient by a discrete central subgroup:
\[
1\longrightarrow\{\pm I\}
\longrightarrow SU(2)
\longrightarrow SO(3)
\longrightarrow1.
\]

This is a model for many Lie-group covering constructions.

\section{Bridge to covering graphs}

The \(SU(2)\to SO(3)\) example shows how a simply connected total space carries
a deck action realizing the base fundamental group.  VIII/14 develops the
same principle combinatorially for graphs, where universal covers become trees
and fundamental groups become free groups.


\section{Exercises}

\begin{exercise}\label{exr:viii13-01}
Define \(SU(2)\).
\end{exercise}
\begin{hint}
Unitary matrices with determinant one.
\end{hint}
\begin{solution}
\(SU(2)=\{A:A^*A=I,\det A=1\}\).
\end{solution}

\begin{exercise}\label{exr:viii13-02}
Write the standard form of an \(SU(2)\) matrix.
\end{exercise}
\begin{hint}
Use two complex parameters.
\end{hint}
\begin{solution}
\(\begin{pmatrix}\alpha&-\overline\beta\\ \beta&\overline\alpha\end{pmatrix}\) with \(|\alpha|^2+|\beta|^2=1\).
\end{solution}

\begin{exercise}\label{exr:viii13-03}
Why is \(SU(2)\cong S^3\)?
\end{exercise}
\begin{hint}
Count real coordinates.
\end{hint}
\begin{solution}
\((\alpha,\beta)\in\mathbb C^2\cong\mathbb R^4\) lies on the unit sphere.
\end{solution}

\begin{exercise}\label{exr:viii13-04}
What group is another model of \(SU(2)\)?
\end{exercise}
\begin{hint}
Use quaternions.
\end{hint}
\begin{solution}
The unit quaternions.
\end{solution}

\begin{exercise}\label{exr:viii13-05}
Which quaternion subspace is \(\mathbb R^3\)?
\end{exercise}
\begin{hint}
Discard the real part.
\end{hint}
\begin{solution}
\(\operatorname{Im}\mathbb H\).
\end{solution}

\begin{exercise}\label{exr:viii13-06}
Define \(R_q\).
\end{exercise}
\begin{hint}
Conjugate imaginary quaternions.
\end{hint}
\begin{solution}
\(R_q(v)=qvq^{-1}\).
\end{solution}

\begin{exercise}\label{exr:viii13-07}
In which group does \(R_q\) lie?
\end{exercise}
\begin{hint}
It preserves orientation and norm.
\end{hint}
\begin{solution}
\(SO(3)\).
\end{solution}

\begin{exercise}\label{exr:viii13-08}
What is \(\ker\rho\)?
\end{exercise}
\begin{hint}
Find unit quaternions commuting with all imaginary ones.
\end{hint}
\begin{solution}
\(\{\pm1\}\).
\end{solution}

\begin{exercise}\label{exr:viii13-09}
Why do \(q\) and \(-q\) give the same rotation?
\end{exercise}
\begin{hint}
The signs cancel in conjugation.
\end{hint}
\begin{solution}
\((-q)v(-q)^{-1}=qvq^{-1}\).
\end{solution}

\begin{exercise}\label{exr:viii13-10}
Write the quaternion for rotation by \(\theta\) about unit axis \(u\).
\end{exercise}
\begin{hint}
Use half-angle.
\end{hint}
\begin{solution}
\(q=\cos(\theta/2)+u\sin(\theta/2)\).
\end{solution}

\begin{exercise}\label{exr:viii13-11}
Why is \(\rho\) surjective?
\end{exercise}
\begin{hint}
Every 3D rotation has axis-angle form.
\end{hint}
\begin{solution}
Choose the half-angle quaternion for that axis and angle.
\end{solution}

\begin{exercise}\label{exr:viii13-12}
How many points are in each fiber of \(\rho\)?
\end{exercise}
\begin{hint}
Use the kernel.
\end{hint}
\begin{solution}
Two.
\end{solution}

\begin{exercise}\label{exr:viii13-13}
What is the deck group?
\end{exercise}
\begin{hint}
Use multiplication by kernel elements.
\end{hint}
\begin{solution}
\(\mathbb Z/2\).
\end{solution}

\begin{exercise}\label{exr:viii13-14}
Identify \(SO(3)\) as a quotient of \(S^3\).
\end{exercise}
\begin{hint}
Identify antipodes.
\end{hint}
\begin{solution}
\(S^3/\{\pm1\}\).
\end{solution}

\begin{exercise}\label{exr:viii13-15}
Identify that quotient as a familiar space.
\end{exercise}
\begin{hint}
Use real projective space.
\end{hint}
\begin{solution}
\(\mathbb RP^3\).
\end{solution}

\begin{exercise}\label{exr:viii13-16}
What is \(\pi_1(SO(3))\)?
\end{exercise}
\begin{hint}
The cover is universal.
\end{hint}
\begin{solution}
\(\mathbb Z/2\).
\end{solution}

\begin{exercise}\label{exr:viii13-17}
What does a \(2\pi\)-rotation loop lift to?
\end{exercise}
\begin{hint}
Start at \(1\).
\end{hint}
\begin{solution}
A path ending at \(-1\).
\end{solution}

\begin{exercise}\label{exr:viii13-18}
Is the \(2\pi\)-rotation loop null-homotopic?
\end{exercise}
\begin{hint}
Its lift is not closed.
\end{hint}
\begin{solution}
No.
\end{solution}

\begin{exercise}\label{exr:viii13-19}
What does a \(4\pi\)-rotation loop lift to?
\end{exercise}
\begin{hint}
Traverse the nontrivial loop twice.
\end{hint}
\begin{solution}
A loop from \(1\) back to \(1\).
\end{solution}

\begin{exercise}\label{exr:viii13-20}
What is the homotopy class of the \(4\pi\)-rotation loop?
\end{exercise}
\begin{hint}
Use the order-two fundamental group.
\end{hint}
\begin{solution}
It is trivial.
\end{solution}

\begin{exercise}\label{exr:viii13-21}
What representation gives the matrix version of \(\rho\)?
\end{exercise}
\begin{hint}
Use conjugation on the Lie algebra.
\end{hint}
\begin{solution}
The adjoint representation \(\operatorname{Ad}:SU(2)\to SO(3)\).
\end{solution}

\begin{exercise}\label{exr:viii13-22}
Write the short exact sequence.
\end{exercise}
\begin{hint}
Kernel, source, quotient.
\end{hint}
\begin{solution}
\(1\to\{\pm I\}\to SU(2)\to SO(3)\to1\).
\end{solution}

\begin{exercise}\label{exr:viii13-23}
Why is the kernel central?
\end{exercise}
\begin{hint}
\(\pm I\) commute with every matrix.
\end{hint}
\begin{solution}
Both scalar matrices lie in the center of \(SU(2)\).
\end{solution}

\begin{exercise}\label{exr:viii13-24}
What is the bridge to VIII/14?
\end{exercise}
\begin{hint}
Compare universal covers.
\end{hint}
\begin{solution}
Graphs have simply connected universal covers that are trees, with free groups acting as deck groups.
\end{solution}


\section{Solved problem dossiers}

\begin{problem}[Three-sphere parametrization]\label{prob:viii13-01}
Verify the usual two-complex-parameter form of \(SU(2)\).
\end{problem}
\begin{solution}
Unitarity forces the second column to be a unit orthogonal vector to the first, hence \((-\overline\beta,\overline\alpha)^T\) up to phase; determinant one fixes the phase. The norm condition is \(|\alpha|^2+|\beta|^2=1\).
\end{solution}

\begin{problem}[Quaternion norm preservation]\label{prob:viii13-02}
Why does \(v\mapsto qvq^{-1}\) preserve Euclidean norm on imaginary quaternions?
\end{problem}
\begin{solution}
Quaternion norm is multiplicative, so \(|qvq^{-1}|=|q||v||q|^{-1}=|v|\). Conjugation also preserves the imaginary subspace.
\end{solution}

\begin{problem}[Orientation]\label{prob:viii13-03}
Why does quaternion conjugation land in \(SO(3)\), not merely \(O(3)\)?
\end{problem}
\begin{solution}
The map depends continuously on \(q\in S^3\), and at \(q=1\) it has determinant \(1\). Since determinant takes values \(\pm1\) and \(S^3\) is connected, the determinant is always \(1\).
\end{solution}

\begin{problem}[Kernel computation]\label{prob:viii13-04}
Prove the kernel is exactly \(\{\pm1\}\).
\end{problem}
\begin{solution}
An element in the kernel commutes with \(i,j,k\), hence with all quaternions. The center of \(\mathbb H\) is \(\mathbb R\); a real unit quaternion is \(\pm1\).
\end{solution}

\begin{problem}[Half-angle check]\label{prob:viii13-05}
Why does a \(2\pi\) rotation correspond to the quaternion \(-1\)?
\end{problem}
\begin{solution}
Set \(\theta=2\pi\): \(q=\cos\pi+u\sin\pi=-1\).
\end{solution}

\begin{problem}[Double cover]\label{prob:viii13-06}
Why does each rotation have exactly two quaternion lifts?
\end{problem}
\begin{solution}
Surjectivity gives at least one lift \(q\). The fiber is the coset \(q\ker\rho=\{q,-q\}\).
\end{solution}

\begin{problem}[Projective identification]\label{prob:viii13-07}
Derive \(SO(3)\cong\mathbb RP^3\).
\end{problem}
\begin{solution}
Use \(SU(2)\cong S^3\) and quotient by the central kernel \(\{\pm1\}\), which is exactly the antipodal equivalence relation defining \(\mathbb RP^3\).
\end{solution}

\begin{problem}[Universal cover]\label{prob:viii13-08}
Why is \(SU(2)\to SO(3)\) universal?
\end{problem}
\begin{solution}
\(SU(2)\cong S^3\) and \(S^3\) is simply connected. A covering with simply connected total space is universal.
\end{solution}

\begin{problem}[Fundamental group]\label{prob:viii13-09}
Recover \(\pi_1(SO(3))\) from deck transformations.
\end{problem}
\begin{solution}
For a universal cover, the deck group is isomorphic to the base fundamental group. The deck group here is \(\{\pm I\}\cong\mathbb Z/2\).
\end{solution}

\begin{problem}[Two-pi loop]\label{prob:viii13-10}
Construct a loop in \(SO(3)\) representing the nontrivial element.
\end{problem}
\begin{solution}
Fix an axis \(u\) and project \(q(t)=\cos(\pi t)+u\sin(\pi t)\). The endpoints \(1,-1\) project to the same rotation, so the projection is a loop whose lift is not closed.
\end{solution}

\begin{problem}[Four-pi loop]\label{prob:viii13-11}
Explain why a \(4\pi\) rotation path is null-homotopic.
\end{problem}
\begin{solution}
It represents twice the generator of \(\mathbb Z/2\), hence the identity. Equivalently its lift to \(S^3\) is closed, and \(S^3\) is simply connected.
\end{solution}

\begin{problem}[Fiber endpoint test]\label{prob:viii13-12}
How can one decide which element of \(\pi_1(SO(3))\) a loop represents?
\end{problem}
\begin{solution}
Lift it from \(I\in SU(2)\). If the lift ends at \(I\), the class is trivial; if it ends at \(-I\), it is the nontrivial class.
\end{solution}

\begin{problem}[Axis reversal]\label{prob:viii13-13}
Compare axis-angle data \((u,\theta)\) and \((-u,-\theta)\).
\end{problem}
\begin{solution}
They give the same quaternion \(\cos(\theta/2)+u\sin(\theta/2)\), hence the same rotation.
\end{solution}

\begin{problem}[Angle \(2\pi-\theta\)]\label{prob:viii13-14}
Explain why distinct quaternion signs can encode the same rotation.
\end{problem}
\begin{solution}
Changing \(q\) to \(-q\) alters the quaternion representative but not conjugation. Axis-angle parametrizations reflect this double-valuedness at the group level.
\end{solution}

\begin{problem}[Central quotient]\label{prob:viii13-15}
Why is \(SU(2)/\{\pm I\}\) naturally a group?
\end{problem}
\begin{solution}
\(\{\pm I\}\) is a normal central subgroup, so the quotient group is defined. The homomorphism theorem identifies it with \(SO(3)\).
\end{solution}

\begin{problem}[Deck action]\label{prob:viii13-16}
Describe the nontrivial deck transformation explicitly.
\end{problem}
\begin{solution}
It is \(q\mapsto -q\), or matrix multiplication \(A\mapsto -A\). It preserves every \(SO(3)\) image.
\end{solution}

\begin{problem}[Quaternion example]\label{prob:viii13-17}
What rotation is represented by \(q=\cos(\pi/4)+k\sin(\pi/4)\)?
\end{problem}
\begin{solution}
The half-angle is \(\pi/4\), so the rotation angle is \(\pi/2\) about the \(k\)-axis, i.e. the \(z\)-axis.
\end{solution}

\begin{problem}[Quaternion \(i\)]\label{prob:viii13-18}
What rotation does \(q=i\) represent?
\end{problem}
\begin{solution}
Write \(i=\cos(\pi/2)+i\sin(\pi/2)\); it represents rotation by \(\pi\) about the \(x\)-axis.
\end{solution}

\begin{problem}[Adjoint viewpoint]\label{prob:viii13-19}
Why does the adjoint action have three real dimensions?
\end{problem}
\begin{solution}
The Lie algebra \(\mathfrak{su}(2)\) is three-dimensional over \(\mathbb R\), and its invariant inner product identifies the adjoint action with orthogonal rotations of \(\mathbb R^3\).
\end{solution}

\begin{problem}[Bridge to graph covers]\label{prob:viii13-20}
State the common pattern between \(SU(2)\to SO(3)\) and a graph universal cover.
\end{problem}
\begin{solution}
A simply connected total space carries a free deck action by the base fundamental group, and the quotient by that action recovers the base.
\end{solution}
